% !TEX root = ../main.tex

\enlargethispage{6\baselineskip}

この部では、本書で扱った圏論と Lean の基礎を、より発展的な学習へ接続します。
ここまでの本文では、ソフトウェアエンジニアが実務で使いやすい概念に絞ってきました。
圏、同型、積、余積、関手、自然変換、モノイド、モナド、Pullback などは、型、関数、データ変換、マイグレーション、API adapter、DB join と結びつけやすい概念です。
しかし圏論には、さらに強力で抽象的な概念が多く存在します。

この部で最初に扱う発展概念は、随伴です。
随伴は、自由構成と忘却、抽象化と具体化、最小の構造を自動的に作る操作などを理解するための中心的な概念です。
たとえば、リストを自由モノイドとして見ると、「要素を並べる」という単純な構造が、どのように普遍的な性質を持つかが見えてきます。

次に、米田の補題を直感的に紹介します。
米田の補題は圏論の中でも特に重要ですが、初学者にとっては抽象度が高い補題です。
本書では完全な証明には踏み込まず、「値や対象は、それに対するすべての観測方法によって特徴づけられる」という直感に絞ります。
この見方は、API 設計、テストダブル、観測可能性、インターフェース設計への比喩として有用です。

さらに、トポス、型理論、仕様言語への橋渡しを行います。
トポスは「集合のように振る舞う世界」として理解でき、論理や幾何、型理論と深くつながっています。
ただし、本書の主目的は実務入門であるため、トポス理論そのものを詳細には展開しません。
ここでは、さらに学ぶための地図を示すに留めます。

最後に、Mathlib の \texttt{CategoryTheory} へ進む道筋を示します。
本書では、理解を優先して、関数同型の \texttt{Iso}、\texttt{Option} と \texttt{List} に限定した \texttt{NatTransOptionList}、\texttt{List.map} の関手則などを個別の小さなモデルとして扱いました。
一般の圏、関手、自然変換へ進むときは、Mathlib の既存ライブラリを使います。
これらの具体例が Mathlib の定義にどのように対応するかを確認し、より高度な Lean 形式化へ進む準備をします。

\section*{この部で扱うこと}

\begin{itemize}
\item 随伴の直感と自由構成
\item Galois接続との関係
\item 米田の補題の直感
\item トポスと内部言語への導入
\item Mathlib の \texttt{CategoryTheory} への橋渡し
\item 次に読むべき教材の位置づけ
\end{itemize}

\section*{この部で扱わないこと}

\begin{itemize}
\item 米田の補題の完全な形式証明
\item Kan拡張の詳細
\item トポス理論の本格展開
\item \(\infty\)圏論の技術的詳細
\item Mathlib CategoryTheory API の網羅的解説
\end{itemize}

\section*{本書全体の締めくくり}

本書の目的は、読者を圏論研究者や形式手法の専門家にすることではありません。
目的は、変換、合成、同値性、保存性という考え方を使って、ソフトウェアの仕様や変更をより正確に扱えるようにすることです。
圏論は、そのための整理された言語を与えます。
Lean 4 は、その言語で書いた主張を機械的に確認する場を与えます。

ここから先は、読者自身の関心に応じて道が分かれます。
実務で Lean を使いたいなら、仕様証明とマイグレーション証明を小さく始めるとよいでしょう。
圏論を深めたいなら、Basic Category Theory、Category Theory in Context、Seven Sketches in Compositionality などへ進むとよいでしょう。
Lean の数学ライブラリを使いたいなら、Mathematics in Lean と Mathlib のドキュメントに進むとよいでしょう。
